\section{Executive function in ADHD}\label{sec:adhd}
ADHD is a heterogeneous disorder in which executive difficulties are common but not universal. A meta-analysis of 83 studies found that groups with ADHD underperform controls on every executive task class examined, most strongly on response inhibition, vigilance, working memory and planning, with medium effect sizes ($0.46$--$0.69$) \citep{willcutt2005ef}. The same authors caution that executive weaknesses are neither necessary nor sufficient to cause all cases of ADHD, and that fewer than half of the (mostly child) samples show impairment on any single executive test \citep{willcutt2005ef}. Any argument that treats ADHD as uniformly executive-impaired overstates the evidence.

Even so, in the subgroup with executive weaknesses, the conjunction is troubling for restructuring-heavy therapy: baseline executive capacity is low, anxiety lowers it further, and, per Attentional Control Theory, the compensatory effort that anxious people deploy is itself an executive resource \citep{eysenck2007act}. The comorbidity of anxiety and ADHD in adults makes this a clinically common situation.
\citegap{prevalence of anxiety in adult ADHD; anxiety-disorder CBT outcome when ADHD is comorbid}
